\chapter{General Properties of Schemes}

\section{Prime Spectra}

\begin{definition}
Let \(A\) be a ring.  The \emph{prime spectrum} \(\Spec A\) is the set of prime
ideals of \(A\).  For an ideal \(I\subset A\), put
\[
    V(I)=\{\mathfrak p\in\Spec A: I\subset \mathfrak p\}.
\]
The closed subsets of \(\Spec A\) are the subsets \(V(I)\).  This topology is
called the \emph{Zariski topology}.
\end{definition}

\begin{proposition}
The subsets \(V(I)\) are the closed subsets of a topology on \(\Spec A\).
\end{proposition}

\begin{proof}
We have \(V(0)=\Spec A\) and \(V(A)=\varnothing\).  If \(I,J\subset A\), then
\[
    V(I)\cup V(J)=V(IJ).
\]
Indeed, a prime ideal contains \(IJ\) if and only if it contains \(I\) or
\(J\).  If \((I_\lambda)\) is any family of ideals, then
\[
    \bigcap_\lambda V(I_\lambda)=V\left(\sum_\lambda I_\lambda\right).
\]
Thus arbitrary intersections and finite unions of the sets \(V(I)\) are again
of the same form.  These are exactly the axioms for closed subsets of a
topology.
\end{proof}

\begin{definition}
For \(f\in A\), the \emph{basic open subset} \(D(f)\subset \Spec A\) is
\[
    D(f)=\{\mathfrak p\in\Spec A:f\notin\mathfrak p\}.
\]
\end{definition}

\begin{proposition}
The basic open subsets \(D(f)\) form a basis for the Zariski topology.  Moreover
\[
    D(f)\cap D(g)=D(fg).
\]
\end{proposition}

\begin{proof}
Since \(D(f)=\Spec A\setminus V((f))\), every \(D(f)\) is open, and
\(D(f)\cap D(g)=D(fg)\) follows from the definition.  If \(U\subset\Spec A\) is
open and \(\mathfrak p\in U\), then \(\Spec A\setminus U=V(I)\) for some ideal
\(I\) after replacing \(I\) by the ideal of functions vanishing on the closed
set.  Since \(\mathfrak p\notin V(I)\), some \(f\in I\) is not in
\(\mathfrak p\).  Then \(\mathfrak p\in D(f)\), and \(D(f)\cap V(I)=\varnothing\),
so \(D(f)\subset U\).  Thus every open set is a union of basic opens.
\end{proof}

\begin{proposition}
For \(f\in A\), the map
\[
    \Spec A_f\longrightarrow D(f),\qquad
    \mathfrak q\longmapsto \{a\in A:a/1\in\mathfrak q\},
\]
is a homeomorphism.
\end{proposition}

\begin{proof}
Prime ideals of \(A_f\) correspond to prime ideals of \(A\) not containing
\(f\), by contraction and extension.  Under this bijection, the basic open
\(D(a/f^n)\subset\Spec A_f\) corresponds to \(D(a)\cap D(f)\subset D(f)\).
Therefore the bijection identifies bases of open subsets, hence is a
homeomorphism.
\end{proof}

\section{The Structure Sheaf}

\begin{definition}
Let \(A\) be a ring and \(X=\Spec A\).  Define \(\OO_X(U)\) to be the set of
functions
\[
    s:U\longrightarrow \coprod_{\mathfrak p\in U} A_{\mathfrak p}
\]
such that \(s(\mathfrak p)\in A_{\mathfrak p}\) and locally \(s\) is represented
by a fraction: for every \(\mathfrak p\in U\), there are \(a,f\in A\) and an
open neighbourhood \(V\subset U\) of \(\mathfrak p\) such that \(f\notin
\mathfrak q\) and \(s(\mathfrak q)=a/f\in A_{\mathfrak q}\) for all
\(\mathfrak q\in V\).  With pointwise operations this is a sheaf of rings,
called the \emph{structure sheaf}.
\end{definition}

\begin{proposition}
Let \(X=\Spec A\).  For every \(f\in A\), the natural map
\[
    A_f\longrightarrow \OO_X(D(f))
\]
is an isomorphism.  In particular \(\OO_X(X)\simeq A\).
\end{proposition}

\begin{proof}
The map sends \(a/f^n\) to the section \(\mathfrak p\mapsto a/f^n\) on
\(D(f)\).  It is injective: if \(a/f^n\) maps to zero, then for every prime
\(\mathfrak p\not\supset f\) the element \(a/1\) is zero in \(A_{\mathfrak p}\).
Hence the support of \(a\) in \(\Spec A_f\) is empty.  Equivalently,
\((a/1)=0\) in \(A_f\), so \(a/f^n=0\).

For surjectivity, let \(s\in\OO_X(D(f))\).  The definition gives a cover of
\(D(f)\) by basic opens \(D(g_i)\subset D(f)\) such that
\[
    s|_{D(g_i)}=a_i/g_i^{n_i}
    \quad\text{in } A_{g_i}.
\]
It is enough to treat a finite subcover.  Such a finite subcover exists because
\(D(f)\simeq\Spec A_f\) is quasi-compact: if \(D(f)\subset\bigcup_iD(g_i)\),
then \(V((g_i))\cap D(f)=\varnothing\), so \(f\) lies in the radical of the
ideal generated by finitely many \(g_i\)'s.  On overlaps the fractions
\(a_i/g_i^{n_i}\) agree.  After replacing \(g_i\) by powers, write
\(s=a_i/g_i\) on \(D(g_i)\), and choose \(N\) with
\[
    f^N=\sum_i b_i g_i
\]
in \(A\).  The agreement on overlaps implies that, after multiplying by a
large power of \(f\), the elements \(g_j a_i\) and \(g_i a_j\) agree in
\(A_f\).  Then
\[
    a=\sum_i b_i a_i\in A_f
\]
defines an element \(a/f^N\in A_f\) whose restriction to every \(D(g_i)\) is
\(s\).  Since sections of a sheaf are determined locally, \(a/f^N\) maps to
\(s\).
\end{proof}

\begin{corollary}
For \(X=\Spec A\) and \(\mathfrak p\in X\), the stalk \(\OO_{X,\mathfrak p}\)
is canonically isomorphic to the local ring \(A_{\mathfrak p}\).
\end{corollary}

\begin{proof}
The basic open neighbourhoods \(D(f)\) with \(f\notin\mathfrak p\) form a basis
at \(\mathfrak p\).  Therefore
\[
    \OO_{X,\mathfrak p}
    =\varinjlim_{\mathfrak p\in D(f)}\OO_X(D(f))
    \simeq \varinjlim_{f\notin\mathfrak p} A_f
    \simeq A_{\mathfrak p}.
\]
\end{proof}

\section{Locally Ringed Spaces and Schemes}

\begin{definition}
A \emph{locally ringed space} is a pair \((X,\OO_X)\), where \(X\) is a
topological space and \(\OO_X\) is a sheaf of rings, such that every stalk
\(\OO_{X,x}\) is a local ring.  A morphism of locally ringed spaces
\[
    f:(X,\OO_X)\longrightarrow (Y,\OO_Y)
\]
is a continuous map \(f:X\to Y\) together with a morphism of sheaves of rings
\[
    f^\#:\OO_Y\longrightarrow f_*\OO_X
\]
such that the induced local homomorphism
\[
    f^\#_x:\OO_{Y,f(x)}\longrightarrow \OO_{X,x}
\]
maps the maximal ideal of \(\OO_{Y,f(x)}\) into the maximal ideal of
\(\OO_{X,x}\).
\end{definition}

\begin{definition}
An \emph{affine scheme} is a locally ringed space isomorphic to
\((\Spec A,\OO_{\Spec A})\) for some ring \(A\).  A \emph{scheme} is a locally
ringed space \((X,\OO_X)\) admitting an open cover \(X=\bigcup_i U_i\) such
that each \((U_i,\OO_X|_{U_i})\) is an affine scheme.
\end{definition}

\begin{proposition}
For rings \(A\) and \(B\), there is a natural bijection
\[
    \Hom_{\mathrm{Sch}}(\Spec B,\Spec A)\simeq \Hom_{\mathrm{Ring}}(A,B).
\]
\end{proposition}

\begin{proof}
Let \(\varphi:A\to B\) be a ring homomorphism.  It defines a continuous map
\(f:\Spec B\to\Spec A\) by \(f(\mathfrak q)=\varphi^{-1}(\mathfrak q)\).  On
basic opens, \(f^{-1}D(a)=D(\varphi(a))\).  The maps
\[
    A_a\longrightarrow B_{\varphi(a)}
\]
therefore define compatible maps
\(\OO_{\Spec A}(D(a))\to\OO_{\Spec B}(f^{-1}D(a))\), hence a morphism of
sheaves \(f^\#:\OO_{\Spec A}\to f_*\OO_{\Spec B}\).  On stalks it is the local
homomorphism \(A_{\mathfrak p}\to B_{\mathfrak q}\), so \(f\) is a morphism of
locally ringed spaces.

Conversely, if \(f:\Spec B\to\Spec A\) is a morphism of locally ringed spaces,
then applying global sections gives a ring homomorphism
\[
    A=\OO_{\Spec A}(\Spec A)\longrightarrow
    \OO_{\Spec B}(\Spec B)=B.
\]
The construction of the previous paragraph from this global homomorphism gives
the same continuous map because local homomorphisms pull back prime ideals to
prime ideals, and it gives the same sheaf map because the structure sheaf on
basic opens is obtained by localization from global sections.  The two
constructions are inverse.
\end{proof}

\section{Basic Properties}

\begin{definition}
A scheme \(X\) is \emph{reduced}, \emph{irreducible}, \emph{integral}, or
\emph{locally noetherian} if it has an affine open cover
\(\Spec A_i\) with the corresponding rings reduced, with the topological spaces
irreducible, with the rings domains, or with the rings noetherian,
respectively.  It is \emph{noetherian} if it is locally noetherian and its
underlying topological space is quasi-compact.
\end{definition}

\begin{proposition}
For a ring \(A\), the affine scheme \(\Spec A\) is irreducible if and only if
the nilradical of \(A\) is prime.  It is reduced if and only if \(A\) is
reduced.  It is integral if and only if \(A\) is a domain.
\end{proposition}

\begin{proof}
The closed irreducible subsets of \(\Spec A\) are exactly the sets
\(V(\mathfrak p)\) for prime ideals \(\mathfrak p\).  Since the closure of the
generic point \((0)\) is \(V(0)=\Spec A\), the whole space is irreducible
exactly when \(V(0)=V(\sqrt{0})\) is irreducible, which is equivalent to
\(\sqrt{0}\) being prime.  The scheme is reduced exactly when every local ring
\(A_{\mathfrak p}\) is reduced, and this is equivalent to \(A\) being reduced:
if \(a\) is nilpotent in \(A\), it remains nilpotent in every localization; if
\(a/1=0\) in all \(A_{\mathfrak p}\), then the support of \(a\) is empty, so
\(a=0\).  The final assertion follows by combining the first two: an affine
scheme is integral exactly when it is reduced and irreducible, which means
\(\sqrt{0}=0\) is prime, i.e. \(A\) is a domain.
\end{proof}

\begin{definition}
The \emph{dimension} of a scheme \(X\) is the supremum of the lengths \(n\) of
chains of irreducible closed subsets
\[
    Z_0\subsetneq Z_1\subsetneq \cdots \subsetneq Z_n.
\]
For \(X=\Spec A\), this is the Krull dimension of \(A\).
\end{definition}

\begin{example}
If \(k\) is a field, \(\Spec k\) has one point and is the final affine
\(k\)-scheme.  The affine \(n\)-space over \(k\) is
\[
    \AA^n_k=\Spec k[x_1,\ldots,x_n].
\]
The closed points of \(\AA^n_k\) with residue field \(k\) correspond to
\(n\)-tuples in \(k^n\).
\end{example}

\begin{definition}
The projective \(n\)-space over a ring \(A\) is the scheme obtained by gluing
the affine schemes
\[
    U_i=\Spec A\left[\frac{x_0}{x_i},\ldots,\frac{x_n}{x_i}\right],
    \qquad 0\le i\le n,
\]
along the evident common principal opens where both \(x_i\) and \(x_j\) are
non-zero.  It is denoted \(\PP^n_A\).
\end{definition}

\begin{remark}
Equivalently, \(\PP^n_A=\Proj A[x_0,\ldots,x_n]\).  The gluing description is
sufficient for the arguments in these notes: it gives the standard affine open
cover \(U_i=D_+(x_i)\), and intersections \(U_i\cap U_j\) are affine principal
opens.
\end{remark}

\section{Subschemes and Gluing}

\begin{definition}
Let \(X\) be a scheme.  An \emph{open subscheme} is an open subset
\(U\subset X\) with the restricted sheaf \(\OO_X|_U\).  A \emph{closed
subscheme} of \(X\) is a closed subset \(Z\subset X\) together with a sheaf
\(\OO_Z\) such that, locally on \(X\), it is of the form
\[
    \Spec(A/I)\hookrightarrow \Spec A.
\]
Equivalently, a closed subscheme is determined by a quasi-coherent sheaf of
ideals \(\II\subset\OO_X\), with underlying sheaf of rings
\(\OO_X/\II\) on its support.
\end{definition}

\begin{proposition}
On an affine scheme \(X=\Spec A\), closed subschemes are in bijection with
ideals \(I\subset A\).  The closed subscheme associated to \(I\) is
\(\Spec(A/I)\), and its underlying topological space is \(V(I)\).
\end{proposition}

\begin{proof}
An ideal \(I\subset A\) gives the quotient map \(A\to A/I\), hence a morphism
\(\Spec(A/I)\to\Spec A\).  The image consists of prime ideals containing \(I\),
namely \(V(I)\), and the map is a homeomorphism onto this closed subset because
prime ideals of \(A/I\) are exactly prime ideals of \(A\) containing \(I\).
The map of structure sheaves is locally the quotient
\[
    A_f\longrightarrow (A/I)_f\simeq A_f/IA_f,
\]
so it is surjective.  Conversely, a closed subscheme of \(\Spec A\) determines
the kernel of \(A=\Gamma(X,\OO_X)\to\Gamma(Z,\OO_Z)\), and locally the closed
subscheme is obtained by quotienting by this ideal.  These constructions are
inverse.
\end{proof}

\begin{definition}
A \emph{locally closed subscheme} of \(X\) is a closed subscheme of an open
subscheme of \(X\).  Its underlying topological space is a locally closed
subset of \(X\).
\end{definition}

\begin{proposition}[Gluing schemes]
Let \(\{X_i\}\) be schemes.  For each pair \(i,j\), let
\[
    U_{ij}\subset X_i,\qquad U_{ji}\subset X_j
\]
be open subschemes, and suppose we are given isomorphisms
\(\varphi_{ij}:U_{ji}\to U_{ij}\) satisfying
\(\varphi_{ii}=\id\), \(\varphi_{ij}=\varphi_{ji}^{-1}\), and the cocycle
condition on triple overlaps.  Then there exists a scheme \(X\), unique up to
unique isomorphism, obtained by gluing the \(X_i\) along the \(\varphi_{ij}\).
\end{proposition}

\begin{proof}
First glue the underlying topological spaces by taking the quotient of
\(\coprod_i X_i\) under the equivalence relation generated by
\(\varphi_{ij}\).  The hypotheses ensure that each \(X_i\) maps homeomorphically
onto an open subset of the quotient and that these open subsets cover the
quotient.  Glue the sheaves \(\OO_{X_i}\) by the sheaf gluing proposition of
the first chapter, using the isomorphisms \(\varphi_{ij}^\#\) on overlaps.
Every point has an open neighbourhood contained in some \(X_i\), hence an
affine neighbourhood.  Thus the glued locally ringed space is a scheme.
Uniqueness follows because any scheme glued from the same data has the same
open cover and the same transition isomorphisms.
\end{proof}

\section{The Proj Construction}

\begin{definition}
Let \(S=\bigoplus_{n\ge0}S_n\) be a graded ring.  A homogeneous prime ideal
\(\mathfrak p\subset S\) is \emph{relevant} if it does not contain the
irrelevant ideal \(S_+=\bigoplus_{n>0}S_n\).  The set of relevant homogeneous
prime ideals is denoted \(\Proj S\).  For a homogeneous ideal \(I\), put
\[
    V_+(I)=\{\mathfrak p\in\Proj S:I\subset\mathfrak p\}.
\]
These sets are the closed subsets of the Zariski topology on \(\Proj S\).
\end{definition}

\begin{definition}
For a homogeneous element \(f\in S_+\), set
\[
    D_+(f)=\{\mathfrak p\in\Proj S:f\notin\mathfrak p\}.
\]
Let \(S_{(f)}\) be the degree-zero part of the localized graded ring \(S_f\).
\end{definition}

\begin{proposition}
The subsets \(D_+(f)\), for homogeneous \(f\in S_+\), form a basis of
\(\Proj S\), and there is a canonical homeomorphism
\[
    D_+(f)\simeq \Spec S_{(f)}.
\]
\end{proposition}

\begin{proof}
The \(D_+(f)\) are complements of \(V_+((f))\), hence open, and
\[
    D_+(f)\cap D_+(g)=D_+(fg).
\]
Every open subset is a union of such opens because closed subsets are defined
by homogeneous ideals.  A relevant homogeneous prime \(\mathfrak p\) not
containing \(f\) gives a prime ideal \((\mathfrak pS_f)_0\subset S_{(f)}\).
Conversely, a prime ideal \(\mathfrak q\subset S_{(f)}\) gives the homogeneous
prime ideal of elements \(a\in S\) for which \(a/f^n\) has degree zero after a
suitable shift and lies in \(\mathfrak q\).  These maps are inverse and identify
basic opens, hence give a homeomorphism.
\end{proof}

\begin{proposition}
There is a unique structure sheaf on \(\Proj S\) such that
\[
    \OO_{\Proj S}(D_+(f))\simeq S_{(f)}
\]
for every homogeneous \(f\in S_+\).  With this sheaf, \(\Proj S\) is a scheme.
\end{proposition}

\begin{proof}
The rings \(S_{(f)}\) and the localization maps between them define a sheaf on
the basis of opens \(D_+(f)\).  By the sheaf-on-a-basis theorem, this extends
uniquely to a sheaf on \(\Proj S\).  Since each \(D_+(f)\) is identified with
\(\Spec S_{(f)}\) and the sheaf restricts to the affine structure sheaf there,
these opens are affine schemes.  They cover \(\Proj S\), so \(\Proj S\) is a
scheme.
\end{proof}

\begin{example}
For \(S=A[x_0,\ldots,x_n]\) with the standard grading,
\[
    D_+(x_i)\simeq
    \Spec A[x_0/x_i,\ldots,x_n/x_i],
\]
and the schemes \(D_+(x_i)\) glue to the projective space \(\PP^n_A\) defined
above.
\end{example}

\section{Noetherian, Jacobson, Normal, and Regular Schemes}

\begin{fact}
Let \(A\) be a noetherian ring.  Then every descending chain of closed subsets
of \(\Spec A\) stabilizes, every irreducible closed subset has a unique generic
point, and every open subset of \(\Spec A\) is quasi-compact if and only if it
is a finite union of basic opens.
\end{fact}

\begin{proposition}
If \(X\) is a noetherian scheme, then every closed subset of \(X\) has finitely
many irreducible components, each with a unique generic point.
\end{proposition}

\begin{proof}
Cover \(X\) by finitely many affine noetherian opens.  On each affine open, the
assertion is the noetherian topological fact just stated.  A closed subset of
\(X\) is the union of its intersections with the affine opens, so it has only
finitely many irreducible components.  The generic point of an irreducible
component is obtained on any affine open meeting it, and uniqueness follows
from the fact that two points with the same closure are equal in a sober
noetherian affine scheme.
\end{proof}

\begin{definition}
A scheme \(X\) is \emph{normal} if all local rings \(\OO_{X,x}\) are integrally
closed domains.  It is \emph{regular} if all local rings \(\OO_{X,x}\) are
regular local rings.
\end{definition}

\begin{fact}
Let \(A\) be a noetherian domain.  The following are equivalent:
\begin{enumerate}[label=(\roman*)]
    \item \(A\) is normal;
    \item \(A_{\mathfrak p}\) is normal for every prime ideal \(\mathfrak p\);
    \item \(A_{\mathfrak m}\) is normal for every maximal ideal \(\mathfrak m\).
\end{enumerate}
Every regular local ring is a normal domain.  A one-dimensional noetherian
regular local ring is a discrete valuation ring.
\end{fact}

\begin{proposition}
Normality and regularity are local properties on a scheme.
\end{proposition}

\begin{proof}
Let \(X=\bigcup_i U_i\) be an open cover.  The local ring of \(U_i\) at
\(x\in U_i\) is canonically \(\OO_{X,x}\).  Hence if \(X\) is normal or regular,
so is every \(U_i\).  Conversely, if the \(U_i\) are normal or regular, every
point \(x\in X\) belongs to some \(U_i\), and \(\OO_{X,x}=\OO_{U_i,x}\) has the
required property.
\end{proof}

\begin{definition}
Let \(X\) be a scheme over a field \(k\), and let \(x\in X\) be a point with
residue field \(\kappa(x)\).  The \emph{Zariski tangent space} of \(X\) at
\(x\) is
\[
    T_xX=\Hom_{\kappa(x)}(\mathfrak m_x/\mathfrak m_x^2,\kappa(x)),
\]
where \(\mathfrak m_x\) is the maximal ideal of \(\OO_{X,x}\).
\end{definition}

\begin{proposition}
Let \(X=\Spec k[x_1,\ldots,x_n]/(f_1,\ldots,f_r)\), and let \(x\in X\) be a
\(k\)-rational point.  Then \(T_xX\) is the kernel of the Jacobian matrix
\[
    \left(\frac{\partial f_i}{\partial x_j}(x)\right):
    k^n\longrightarrow k^r.
\]
\end{proposition}

\begin{proof}
A tangent vector at \(x\) is a \(k\)-algebra homomorphism
\[
    k[x_1,\ldots,x_n]/(f_i)\longrightarrow k[\varepsilon]/(\varepsilon^2)
\]
lifting the point \(x\).  Such a homomorphism sends
\[
    x_j\longmapsto x_j(x)+v_j\varepsilon.
\]
The relation \(f_i=0\) is preserved if and only if
\[
    0=f_i(x+v\varepsilon)
    =
    f_i(x)+\sum_j \frac{\partial f_i}{\partial x_j}(x)v_j\varepsilon
\]
in \(k[\varepsilon]/(\varepsilon^2)\).  Since \(f_i(x)=0\), the condition is
that the vector \(v=(v_j)\) lie in the displayed kernel.
\end{proof}

\section{Exercises Used Later}

\begin{exercise}[Nilpotents and reduced closed subschemes]\label{ex:reduced-closed}
Let \(X=\Spec A\).  Show that the closed subset \(V(I)\) has a unique reduced
closed subscheme structure, namely \(\Spec(A/\sqrt I)\).
\end{exercise}

\begin{solution}
Any closed subscheme with underlying set \(V(I)\) is \(\Spec(A/J)\) with
\(\sqrt J=\sqrt I\).  It is reduced exactly when \(A/J\) is reduced, which is
equivalent to \(J\) being radical.  The only radical ideal with radical
\(\sqrt I\) is \(\sqrt I\) itself.  Hence the reduced structure is
\(\Spec(A/\sqrt I)\).
\end{solution}

\begin{exercise}[Affine gluing test]\label{ex:affine-gluing-test}
Let \(X=U\cup V\), where \(U\), \(V\), and \(U\cap V\) are affine open
subschemes.  Show that giving a morphism \(X\to Y\) is equivalent to giving
morphisms \(U\to Y\) and \(V\to Y\) that agree on \(U\cap V\).
\end{exercise}

\begin{solution}
One direction is restriction.  Conversely, the two morphisms give compatible
continuous maps and compatible morphisms of structure sheaves on the open cover
\(\{U,V\}\).  By gluing morphisms of sheaves, the sheaf maps glue, and by
gluing continuous maps the underlying maps glue.  The local-ring condition for
a morphism of locally ringed spaces is checked on \(U\) and \(V\).  Therefore
the glued data define a unique morphism \(X\to Y\).
\end{solution}

\begin{exercise}[Closed points of affine space]\label{ex:closed-points-affine}
Let \(k\) be a field and \(X=\AA^n_k\).  Show that closed points of \(X\)
correspond to finite field extensions \(K/k\) generated by \(n\) elements,
together with an \(n\)-tuple in \(K^n\), up to the evident equivalence.
\end{exercise}

\begin{solution}
A closed point is a maximal ideal \(\mathfrak m\subset k[x_1,\ldots,x_n]\).
The residue field \(K=k[x_1,\ldots,x_n]/\mathfrak m\) is a finite extension of
\(k\) by Zariski's lemma.  The images of the variables give an \(n\)-tuple in
\(K^n\) generating \(K\) over \(k\).  Conversely, such an \(n\)-tuple defines a
surjective homomorphism \(k[x_1,\ldots,x_n]\to K\), whose kernel is maximal.
Changing the identification of \(K\) gives the evident equivalence.
\end{solution}

\begin{exercise}[The standard cover of projective space]\label{ex:standard-cover}
Show that the standard opens \(D_+(x_i)\) cover \(\PP^n_A\), and compute
\[
    D_+(x_i)\cap D_+(x_j).
\]
\end{exercise}

\begin{solution}
No relevant homogeneous prime contains all variables \(x_0,\ldots,x_n\), so
the \(D_+(x_i)\) cover.  Their intersection is \(D_+(x_ix_j)\).  Inside
\(D_+(x_i)\simeq\Spec A[x_0/x_i,\ldots,x_n/x_i]\), it is the principal open
where \(x_j/x_i\) is invertible.  This also agrees with the corresponding
principal open inside \(D_+(x_j)\).
\end{solution}

\begin{exercise}[Regular curves have DVR local rings]\label{ex:regular-curve-dvr}
Let \(X\) be a noetherian regular scheme of dimension one.  Show that for every
closed point \(x\in X\), \(\OO_{X,x}\) is a discrete valuation ring.
\end{exercise}

\begin{solution}
The local ring \(\OO_{X,x}\) is noetherian, local, regular, and
one-dimensional.  By the commutative algebra fact above, every
one-dimensional noetherian regular local ring is a DVR.  This fact is used
later to define valuations of rational functions at points of a smooth curve.
\end{solution}

\section{Function Fields, Rational Maps, and Normalization}

\begin{definition}
Let \(X\) be an integral scheme with generic point \(\eta\).  Its
\emph{function field} is
\[
    k(X)=\OO_{X,\eta}.
\]
If \(U=\Spec A\subset X\) is a non-empty affine open, then \(A\) is a domain and
\[
    k(X)=\Frac(A).
\]
\end{definition}

\begin{proposition}
Let \(X\) be integral and let \(U,V\subset X\) be non-empty open subsets.  Then
\[
    \OO_X(U)\subset k(X),\qquad \OO_X(V)\subset k(X)
\]
and the two inclusions agree on \(U\cap V\).  Thus regular functions on
non-empty opens may be regarded as rational functions.
\end{proposition}

\begin{proof}
For \(s\in\OO_X(U)\), its germ at the generic point \(\eta\in U\) is an element
of \(k(X)\).  If this germ is zero, then the section is zero on some
neighbourhood of \(\eta\).  Since \(X\) is irreducible, every non-empty open
subset of \(U\) is dense, and the zero locus of \(s\) is open; by the sheaf
property the section is zero wherever the germ is zero after restricting to an
affine cover.  Hence the map \(\OO_X(U)\to k(X)\) is injective.  Compatibility
on intersections follows because both maps are given by taking the germ at
\(\eta\).
\end{proof}

\begin{definition}
Let \(X,Y\) be schemes with \(X\) reduced.  A \emph{rational map}
\[
    X\dashrightarrow Y
\]
is an equivalence class of morphisms \(U\to Y\), where \(U\subset X\) is a
dense open subset; two such morphisms are equivalent if they agree on a dense
open subset of the intersection of their domains.  If \(Y\) is separated, it is
enough to require agreement on \(U\cap V\).
\end{definition}

\begin{proposition}
Let \(X\) be integral and let \(Y\) be separated.  Two morphisms
\[
    f,g:X\to Y
\]
which agree on a non-empty open subset of \(X\) are equal.
\end{proposition}

\begin{proof}
The pair \((f,g)\) defines a morphism
\[
    X\to Y\times Y.
\]
The locus where \(f=g\) is the inverse image of the diagonal
\(\Delta_Y\subset Y\times Y\).  Since \(Y\) is separated, the diagonal is
closed.  This inverse image is therefore closed in \(X\), and it contains a
non-empty open subset.  In an integral scheme every non-empty open subset is
dense, so the closed subset is all of \(X\).  Hence \(f=g\).
\end{proof}

\begin{definition}
Let \(X\) be an integral scheme with function field \(K=k(X)\), and let
\(L/K\) be an algebraic field extension.  The \emph{normalization of \(X\) in
\(L\)} is the scheme obtained by gluing the affine schemes
\[
    \Spec B_U,
\]
where \(U=\Spec A\subset X\) ranges over affine opens and \(B_U\) is the
integral closure of \(A\) in \(L\).
\end{definition}

\begin{proposition}
The normalization construction glues to a normal integral scheme
\(\nu:X^\nu\to X\).  If \(X\) is noetherian and \(L/K\) is finite, and if the
integral closure of every affine coordinate ring in \(L\) is finite, then
\(\nu\) is finite.
\end{proposition}

\begin{proof}
If \(V=D(f)\subset U=\Spec A\), then the coordinate ring of \(V\) is \(A_f\).
The integral closure of \(A_f\) in \(L\) is \((B_U)_f\): integrality commutes
with localization because a monic equation over \(A_f\) can be cleared of
denominators after multiplying the element by a power of \(f\), and conversely
localization preserves monic equations.  Hence the affine normalizations agree
on overlaps and glue.  The local rings of the glued scheme are integrally
closed domains, so it is normal and integral.  Under the stated finiteness
hypothesis, each affine inverse image is \(\Spec B_U\) with \(B_U\) finite over
\(A\), hence \(\nu\) is finite.
\end{proof}

\begin{fact}
If \(A\) is an excellent noetherian domain, then its integral closure in every
finite field extension of \(\Frac(A)\) is finite over \(A\).  In particular
this holds for finitely generated algebras over a field and for complete
noetherian local domains.
\end{fact}

\begin{corollary}
Let \(X\) be an integral curve of finite type over a field \(k\).  Then the
normalization \(X^\nu\to X\) is finite.
\end{corollary}

\begin{proof}
The curve is covered by affines \(\Spec A\), where \(A\) is a finitely generated
domain over \(k\).  Such rings are excellent, so their integral closures in the
function field are finite.  The previous proposition gives finiteness of the
normalization.
\end{proof}
